\documentclass[12pt]{article}
\usepackage{amsmath,amssymb,booktabs,longtable,array}
\usepackage{fontspec}
\usepackage{polyglossia}
\setmainlanguage{arabic}
\setotherlanguage{english}
\newfontfamily\arabicfont[Script=Arabic]{Amiri}
\usepackage{geometry}
\geometry{margin=2.1cm}
\usepackage{hyperref}

\title{اختبار قانون Scaling للأوزان المحلية متعددة المقاييس}
\author{مشروع تجربة البحث عن الأعداد الأولية}
\date{2026-09-01}

\begin{document}
\maketitle

\section{السؤال}

في التجربة السابقة مثلنا pair kernel بواسطة موجات Haar محلية متعددة المقاييس.
السؤال هنا:

\begin{quote}
هل الأوزان التي نحتاجها تتبع قانونًا يعتمد أساسًا على النسبة
\[
L/H,
\]
بحيث إذا تغير موضع التجربة \(X\) وطول الفترة \(H\)، تنتقل البنية نفسها
بمجرد إعادة تحجيم طول الدعم \(L\)؟
\end{quote}

لو حدث ذلك، فسيكون لدينا مرشح لبنية scale-covariant بدل قائمة أوزان
مرتبطة بنطاق عددي واحد.

\section{تنبيه منهجي: الوزن الخام ليس قابلًا للتحديد دائمًا}

أول نتيجة في التجربة كانت سلبية لكنها مهمة.

إذا كانت الفترات متساوية الطول \(H\)، فإن كثيرًا من موجات Haar ذات:
\[
L<H
\]
تنتج تقريبًا نفس \emph{شكل} covariance بعد التطبيع.

في مثال \(H=40000\)، أعطت المقاييس:
\[
L/H=0.0512,\ 0.1024,\ 0.2048,\ 0.4096,\ 0.8192
\]
كلها تقريبًا:
\[
\rho_1\approx-\frac12,
\qquad
\rho_{k\ge2}\approx0.
\]

إذن يمكن لـNNLS أن ينقل الوزن من مقياس صغير إلى آخر من غير أن يغير
kernel المرصود تقريبًا.

وعليه:

\[
\boxed{
\text{القيم الفردية }v_s\text{ ليست كائنًا قابلًا للتفسير وحدها.}
}
\]

الكائن الأكثر ثباتًا هو \emph{مساهمة التباين}:
\[
c_s=v_sC_s(0),
\]
ثم جمعها في نطاقات بعدية من \(L/H\).

\section{الطيف البعدي الذي اختبرناه}

استخدمنا النطاقات:

\[
L/H<2,
\]
\[
2\le L/H<4,
\]
\[
4\le L/H<8,
\]
\[
L/H\ge8.
\]

هذه النطاقات الخشنة تقلل مشكلة عدم قابلية تحديد المقياس الدقيق.

\section{تجربة تحكم اصطناعية}

أنشأنا فترات متساوية الطول عند:

\[
X=20,\ 40,\ 60\text{ مليونًا},
\]

مع ثلاث نسب:

\[
X/H=250,\ 500,\ 1000.
\]

أي أننا غيرنا \(X\) و\(H\) مع إبقاء \(X/H\) ثابتًا داخل كل مجموعة.

استخدمنا مكتبة مقاييس لوغاريتمية أكثر كثافة من التجربة السابقة:
أربع نقاط في كل octave، ثم أعدنا بعض اختبارات النقل بثمان وست عشرة
نقطة في كل octave.

\section{الانهيار على متغير } \texorpdfstring{\(L/H\)}{L/H}

في الحالات \(X=20\) و\(40\) مليونًا، حيث يتضاعف \(H\) بالضبط داخل
كل مجموعة، كان الطيف البعدي شبه ثابت.

المتوسط بين الحالتين:

\begin{center}
\begin{tabular}{rrrrr}
\toprule
\(X/H\) & \(L/H<2\) & \(2\le L/H<4\) & \(4\le L/H<8\) & \(L/H\ge8\)\\
\midrule
250  & 0.5558 & 0.2126 & 0.1307 & 0.1009\\
500  & 0.5499 & 0.2155 & 0.1301 & 0.1044\\
1000 & 0.5460 & 0.2202 & 0.1277 & 0.1062\\
\bottomrule
\end{tabular}
\end{center}

وأقصى مسافة \(L^1\) بين طيف \(20\) مليونًا وطيف \(40\) مليونًا
داخل المجموعة نفسها كانت:

\[
\boxed{0.0174}.
\]

إذن عند تغيير \(X,H\) مع تثبيت النسبة، تبقى مساهمات النطاقات البعدية
قريبة جدًا.

\section{اختبار أقوى: انقل الطيف بدل إعادة ملاءمته}

لم نكتف بإعادة NNLS في كل حالة.

أخذنا الطيف الملائم عند:
\[
X=20\text{ مليونًا}
\]
ثم نقلنا كل دعم:
\[
L\longmapsto L\frac{H_{\rm target}}{H_{\rm source}},
\]
من غير إعادة ملاءمة أوزان covariance.

مع مكتبة فيها ثماني نقاط في octave، وعند مضاعفة \(H\):
\[
H_{\rm target}/H_{\rm source}=2,
\]
كان أكبر kernel RMSE عبر النسب الثلاث:

\[
\boxed{3.56\times10^{-4}}.
\]

هذا نقل شبه كامل.

عند:
\[
H_{\rm target}/H_{\rm source}=3,
\]
ارتفع أكبر الخطأ إلى:
\[
0.00424.
\]

لكن عندما كثفنا شبكة المقاييس إلى \(16\) نقطة في octave في حالة
\(X/H=500\)، انخفض خطأ النقل بعامل \(3\) إلى:

\[
\boxed{0.000864}.
\]

إذن جزء كبير من فشل النقل غير الثنائي سببه discretization في مكتبة
المقاييس، لا غياب scaling.

\section{ماذا يعني ذلك؟}

في تجارب التحكم النظرية لدينا دليل قوي على:

\[
\boxed{
C(L,H,X)
\text{ يعتمد بدرجة كبيرة على }
L/H
\text{ و }
X/H.
}
\]

أو بصورة مفهومية:

\[
\boxed{
\text{إذا كبر }H\text{، تتحرك المقاييس المهمة معه.}
}
\]

هذا متوقع رياضيًا من حقيقة أن kernel للفترات المتساوية يمكن كتابته
بدلالة \(X/H\) والمسافة بوحدات \(H\).

لذلك يجب عدم تقديم هذه النتيجة كقانون جديد للأوليات؛ هي تحقق أن
تمثيلنا المحلي يحترم تقريبًا scale covariance الموجودة أصلًا في النظرية.

\section{اختبار على فترات مربعات الأوليات الفعلية}

قسمنا الـ\(450\) فترة إلى أربع كتل متتالية في \(X\).

للفترات الأصلية غير المجمعة، كانت مساهمات:

\[
L/H<2
\]

كما يلي:

\[
0.7564,\quad0.7821,\quad0.7544,\quad0.7217.
\]

أي:

\[
\boxed{
\text{المتوسط}=0.7536,
\qquad
\text{الانحراف المعياري}=0.0248.
}
\]

نحو ثلاثة أرباع مساهمة التباين في هذا التمثيل تقع إذن عند دعم أصغر من
ضعفي طول الفترة النموذجي.

أما:
\[
L/H\ge8
\]
فكان متوسط مساهمته:
\[
0.0523
\]
بانحراف معياري:
\[
0.0136.
\]

هذا أكثر استقرارًا من الأوزان الفردية الخام.

\section{اختبار التجميع}

جمعنا أيضًا كل فترتين متتاليتين، وبالتالي كبر \(H\) تقريبًا.

بقي متوسط مساهمة:
\[
L/H<2
\]
نحو:
\[
0.7181,
\]
مع تشتت أكبر بين اللوحات.

إذن البنية الخشنة تتحرك مع \(H\)، لكنها ليست invariant مثالية عندما
تكون الفترات غير متساوية وتختلف prime gaps بشدة.

\section{نقل طيف لوحة إلى لوحة أخرى}

أخذنا طيف اللوحة الأولى فقط، ثم:

\begin{enumerate}
\item غيرنا مقاييس الدعم بحسب نسبة median \(H\).
\item عدلنا amplitude باستعمال العامل النظري التقريبي
\[
\frac{\log(X_1/H_1)}{\log(X_2/H_2)}.
\]
\item لم نعد ملاءمة شكل الطيف على اللوحات الأخرى.
\end{enumerate}

كانت kernel RMSE عند اللوحات \(2,3,4\):

\[
0.00338,\quad0.00164,\quad0.00610.
\]

المتوسط:
\[
\boxed{0.00371}.
\]

إذن النقل يعمل تقريبًا، لكنه أسوأ من حالات التحكم المتساوية.
اللوحة الرابعة كانت الأصعب.

وهذا منطقي: فترات مربعات الأوليات غير متساوية، و\(H_i\) نفسه يتقلب
بقوة بسبب prime gaps، فلا تختزل هندستها بالكامل في median \(H\).

\section{هل وجدنا قانونًا بسيطًا للأوزان؟}

الجواب يحتاج فصل مستويين.

\subsection{الأوزان الفردية}

لا:

\[
\boxed{
v_s\text{ الفردية ليست قانونًا ثابتًا قابلًا للتحديد من kernel وحده.}
}
\]

السبب هو collinearity بين بعض مقاييس Haar.

\subsection{الطيف البعدي الخشن}

نعم، توجد إشارة قوية إلى scaling:

\[
\boxed{
\text{مساهمة covariance تنتظم بصورة أفضل عندما نكتب المقياس كـ }L/H.
}
\]

وفي التحكم المتساوي الطول، يمكن نقل الطيف بين قيم \(X,H\) المختلفة
بدقة عالية جدًا من غير إعادة ملاءمة كاملة.

لكن هذا إلى حد كبير انعكاس للبنية scale-covariant في kernel النظري نفسه،
وليس بعد قانونًا مستقلًا اكتشفناه في الأوليات.

\section{أهم تصحيح مفاهيمي}

في التقرير السابق قلنا إن ``أوزان المقاييس'' قد تكون الهدف التالي.

هذه التجربة تصحح العبارة.

الهدف الصحيح ليس:
\[
v(L).
\]

بل أقرب إلى:
\[
\boxed{
\text{طيف مساهمة بعدي }
\mathcal W(L/H;X/H),
}
\]

مع الانتباه إلى أن التمثيل الدقيق داخل النطاقات الصغيرة قد لا يكون وحيدًا.

أي أن الكائن الذي يستحق البحث هو \emph{equivalence class} من الآليات
المحلية التي تولد kernel نفسه، لا قائمة واحدة من coefficients.

\section{السؤال التالي}

الاختبار الحاسم التالي هو تغيير نطاق \(X\) نفسه بأكثر من العامل الصغير
المتاح في بياناتنا الحالية.

حاليًا:
\[
X\approx 1.6\times10^7
\text{ إلى }
6.3\times10^7.
\]

نحتاج نطاقات أعلى، ثم نختبر هل:
\[
\mathcal W(L/H;X/H)
\]
تظل ثابتة بعد التحكم في \(X/H\)، أم تبدأ تصحيحات جديدة بالظهور.

إذا استمر collapse على عدة مراتب عشرية من \(X\)، يصبح scaling candidate
أقوى بكثير.

\section{ملفات التجربة}

\begin{itemize}
\item \texttt{local\_multiscale\_scaling\_synthetic\_bands.csv}
\item \texttt{local\_multiscale\_scaling\_synthetic\_transfer.csv}
\item \texttt{local\_multiscale\_scaling\_prime\_panels.csv}
\item \texttt{local\_multiscale\_scaling\_prime\_transfer.csv}
\item \texttt{local\_multiscale\_scaling\_identifiability.csv}
\item \texttt{local\_multiscale\_scaling\_summary.csv}
\item \texttt{analyze\_local\_multiscale\_scaling.py}
\item \texttt{plots/local\_multiscale\_scaling\_synthetic\_bands.svg}
\item \texttt{plots/local\_multiscale\_scaling\_prime\_panels.svg}
\item \texttt{plots/local\_multiscale\_scaling\_transfer\_error.svg}
\end{itemize}

\end{document}
